% !PRINTED 345
\weberpage{345}
%<<CONT
given, and further developed by \gesperrt{Encke}\footnote{Crelle's Journal, Bd.~22 (1841).}, and which lends itself particularly well to a practical execution of the numerical computations.

If, as above, $\alpha, \beta, \gamma \ldots$ are the roots of an equation, and $\alpha$ is the absolutely greatest among them, so that none of the others equals the root $\alpha$ in absolute value, then
\weq{1}{\sqrt[m]{\alpha^m + \beta^m + \gamma^m + \cdots}}
has, as $m$ grows without bound, the value $\alpha$ as its limit; indeed the convergence toward $\alpha$ is even better than for the quotient considered in the preceding section,
\[
\frac{\alpha^m + \beta^m + \gamma^m + \cdots}{\alpha^{m-1} + \beta^{m-1} + \gamma^{m-1} + \cdots},
\]
as one recognises when, according to the generalised binomial theorem, one sets:
\[
\begin{aligned}
\sqrt[m]{\alpha^m + \beta^m + \gamma^m + \cdots} &= \alpha\sqrt[m]{1 + \left(\frac{\beta}{\alpha}\right)^m + \left(\frac{\gamma}{\alpha}\right)^m + \cdots} \\
&= \alpha\left[1 + \frac{1}{m}\left(\frac{\beta}{\alpha}\right)^m + \frac{1}{m}\left(\frac{\gamma}{\alpha}\right)^m + \cdots\right]
\end{aligned}
\]

The absolutely smallest among the roots is obtained by the same principle as the limit of
\weq{2}{\dfrac{1}{\sqrt[m]{\dfrac{1}{\alpha^m} + \dfrac{1}{\beta^m} + \dfrac{1}{\gamma^m} + \cdots}},}
and if $p$ is an arbitrary value, one obtains the root lying nearest to the value $p$ as the limit of
\weq{3}{p + \dfrac{1}{\sqrt[m]{\dfrac{1}{(\alpha-p)^m} + \dfrac{1}{(\beta-p)^m} + \cdots}}.}
If one transforms the equation $f(x)$ by the substitution
\[
y = \frac{1}{x}, \qquad y = \frac{1}{x - p}
\]
then the expressions (2), (3) arise from the expression (1), to which we shall therefore now confine ourselves.
